\SourcePage{338}
\SGAchapter{XIX}{Noether 定理}{Le théorème de Noether}
\begin{center}P. Deligne\end{center}

\SGAsection{1}{定理的陈述}
本讲用现代语言重述 Lefschetz [3] 第108页对下述 Noether \Term{theorem}{定理}（1.2）的证明。

\textbf{1.1。}设 $k$ 是\Term{characteristic}{特征}为 $p$ 的\Term{algebraically-closed-field}{代数闭域}。设 $d\geq1$，$\underline a=(a_i)_{1\leq i\leq d}$ 是一列 $\geq1$ 的整数。$k$ 上\Term{projective-space}{射影空间} $\mathbb P^{n+d}$ 中\Term{dimension}{维数}为 $n$、\Term{multidegree}{多重次数}为 $\underline a$ 的\Term{complete-intersection}{完全交} $V$，如果在 $k$ 的某个\Term{field-extension}{扩域}上，与一个由系数在\Term{prime-field}{素域}上\Term{algebraic-independence}{代数独立}的方程所定义的\Term{complete-intersection}{完全交}\Term{projective-isomorphism}{射影同构}，就称它是\Term{generic}{一般的}。以下除非明确另作说明，均假设 $d\geq1$ 且 $a_i\geq2$。

\begin{SGAresult}{\Term{theorem}{定理}}{1.2}
设 $V\subset\mathbb P^{2+d}$ 是\Term{multidegree}{多重次数}为 $\underline a$ 的\Term{generic}{一般}\Term{complete-intersection}{完全交}\Term{surface}{曲面}。假设 $\underline a\ne(2),(3),(2,2)$。则 $V$ 上的每条（\Term{reduced}{约化}）\Term{curve}{曲线}都是 $V$ 与 $\mathbb P^{2+d}$ 中某个\Term{hypersurface}{超曲面}的\Term{scheme-theoretic-intersection}{概形论交}。
\end{SGAresult}
下面说明，该\Term{theorem}{定理}由其推论
\SourcePage{339}
\begin{equation}
\operatorname{Pic}(\mathbb P^{2+d})=\mathbb Z\xrightarrow{\sim}\operatorname{Pic}(V)
\tag{1.2.1}
\end{equation}
\Term{pushout}{推出}。

如果 $C$ 是 $V$ 上的一条\Term{curve}{曲线}，则（1.2.1）蕴含对某个适当的 $n$ 有 $\mathcal O(C)\simeq\mathcal O(n)$；若 $H$ 是\Term{hyperplane-section}{超平面截面}，则 $n=(\prod a_i)^{-1}(C,H)>0$。只需证明 $|nH|$ 在 $V$ 上诱导的\Term{linear-system}{线性系}是\Term{complete-linear-system}{完全的}。将 $V$ 写成\Term{degree}{次数}分别为 $a_i$ 的\Term{hypersurface}{超曲面} $H_i$ 的交，并令 $V_i=H_1\cap\cdots\cap H_i$。对 $i$ 作归纳即可验证 $|nH|$ 在 $V_i$ 上诱导的\Term{linear-system}{线性系}是\Term{complete-linear-system}{完全的}：在\Term{exact-sequence}{正合列}
\[
H^0(V_i,\mathcal O(n))\longrightarrow H^0(V_{i+1},\mathcal O(n))
\longrightarrow H^1(V_i,\mathcal O(n-a_{i+1}))
\]
中，最后一群为零（FAC）。

\begin{SGAresult}{\Term{theorem}{定理}}{1.3}
设 $V\subset\mathbb P^{2n+d}$ 是偶\Term{dimension}{维数} $2n\ne0$、\Term{multidegree}{多重次数}为 $\underline a$ 的\Term{generic}{一般}\Term{complete-intersection}{完全交}。假设下列条件之一成立：
\begin{enumerate}
\item[(a)] $k$ 的\Term{characteristic}{特征}不是 $2$，且 $(2n;\underline a)$ 不等于 $(2n;2)$、$(2n;2,2)$ 或 $(2;3)$。
\item[(b)] \Term{hodge-number}{Hodge 数} $h^{2n,0}$ 非零，即 $\sum(a_i-1)>2n$。
\end{enumerate}
则 $H^{2n}(V,\mathbb Q_\ell(n))$ 中的每个\Term{algebraic-cohomology-class}{代数上同调类}，都是\Term{linear-section}{线性截面}的类 $\eta^n$ 的有理数倍。
\end{SGAresult}

由1.3推导1.2.1。1.2中的例外情形恰好是 $\sum(d_i-1)\leq2$ 的情形；因而可以对1.2中的\Term{surface}{曲面} $V$ 应用1.3（b）。已知 $\operatorname{Pic}^{\tau}(V)=0$（第 XI 讲1.8），故 $\operatorname{Pic}(V)$ \Term{injectivity}{单射}到 $H^2(V,\mathbb Q_\ell(1))$ 中。由1.3，$\operatorname{Pic}(V)$ 的每个元素都是 $\mathcal O(1)$ 的类 $\eta$ 的有理数倍；再由第 XI 讲1.8即可得出结论。

显然，1.2与1.3是否成立只依赖 $p,n,\underline a$，而不依赖所选的\Term{generic}{一般}\Term{complete-intersection}{完全交}。

\SourcePage{340}
\begin{SGAresult}{\Term{theorem}{定理}}{1.4}
设 $X\subset\mathbb P$ 是 $k$ 上纯\Term{dimension}{维数} $2n+1$ 的\Term{smooth}{光滑}\Term{projective-variety}{射影簇}。设 $\ell\subset\check{\mathbb P}$ 是 $X$ 的\Term{hyperplane-section}{超平面截面}的一个\Term{lefschetz-pencil}{Lefschetz 铅笔}。当 $k$ 的\Term{characteristic}{特征}为 $2$ 时，还假设 $\ell$ 是\Term{generic}{一般}\Term{lefschetz-pencil}{Lefschetz 铅笔}。

假设 $X$ 满足（LV）（第 XVIII 讲5.2.2）。设 $V$ 是\Term{pencil}{铅笔} $\ell$ 中的\Term{generic}{一般}\Term{hyperplane-section}{超平面截面}。则下列两个条件之一成立：
\begin{enumerate}
\item[(a)] \Term{cohomology-group}{上同调群} $H^{2n}(V,\mathbb Q_\ell(n))$ 的\Term{vanishing-cohomology}{消失部分} $\operatorname{Ev}(V,\mathbb Q_\ell(n))$ 中，没有非零\Term{algebraic-cohomology-class}{代数上同调类}。
\item[(b)] \Term{algebraic-cohomology-class}{代数上同调类}生成 $\operatorname{Ev}(V,\mathbb Q_\ell(n))$。
\end{enumerate}
\end{SGAresult}

下面明确1.4中“\Term{generic}{一般}”一词的含义。

（a）设 $G$ 是与 $\mathbb P$ \Term{dual-space}{对偶}的\Term{projective-space}{射影空间} $\check{\mathbb P}$ 中直线的\Term{grassmannian}{Grassmann 簇}。称一个\Term{algebraically-closed-field}{代数闭域}的\Term{spectrum}{谱}到 $G$ 的\Term{morphism}{态射}为 $G$ 的“\Term{geometric-point}{几何点}”。每个\Term{geometric-point}{几何点} $\bar s\to G$ 都定义了将 $X$ 从 $k$ \Term{extension-of-scalars}{扩张标量}到 $k(\bar s)$ 后所得\Term{variety}{簇}的一个\Term{hyperplane-section}{超平面截面}\Term{pencil}{铅笔} $\ell$。若 $\bar s$ 位于 $G$ 的\Term{generic-point}{泛点}之上，就称 $\ell$ 是\Term{generic}{一般的}。

（b）设 $\ell$ 是一个\Term{hyperplane-section}{超平面截面}\Term{pencil}{铅笔}（$\check{\mathbb P}$ 中的一条直线）。每个\Term{geometric-point}{几何点} $\bar s\to\ell$ 都定义了将 $X$ 从 $k$ \Term{extension-of-scalars}{扩张标量}到 $k(\bar s)$ 后所得\Term{variety}{簇}的一个\Term{hyperplane-section}{超平面截面}。若 $\bar s$ 位于 $\ell$ 的\Term{generic-point}{泛点}之上，就称此\Term{hyperplane-section}{超平面截面}是\Term{generic}{一般的}。

取上述 $\bar s$，其在 $\ell$ 中的像为\Term{generic-point}{泛点} $s$；设 $S\subset\ell$ 是对应于\Term{singular-hyperplane-section}{奇异超平面截面}的\Term{finite}{有限}点集。\Term{monodromy-group}{单值群} $\pi_1(\ell-S,\bar s)$ 是 $\operatorname{Gal}(k(\bar s)/k(s))$ 的\Term{quotient-group}{商群}，\Term{action}{作用}在 $\operatorname{Ev}(V,\mathbb Q_\ell(n))$ 上。若 $k$ 的\Term{characteristic}{特征}不是 $2$，这个\Term{action}{作用}是\Term{irreducible}{不可约的}（第 XVIII 讲6.7）；下面第2节将证明，对\Term{generic}{一般}\Term{pencil}{铅笔}，这个结论不需任何\Term{characteristic}{特征}假设仍然成立。$\operatorname{Ev}(V,\mathbb Q_\ell(n))$ 中由\Term{algebraic-cohomology-class}{代数上同调类}生成的\Term{vector-subspace}{子空间}在\Term{monodromy-group}{单值群}（即 Galois 群）\Term{action}{作用}下\Term{stable}{稳定}，因而它或者是全\Term{space}{空间}，或者是零。这就证明了1.4。

\SourcePage{341}
\begin{SGAresult}{推论}{1.5}
设 $V\subset\mathbb P^{2n+d}$ 是偶\Term{dimension}{维数} $2n$、\Term{multidegree}{多重次数}为 $\underline a$ 的\Term{generic}{一般}\Term{complete-intersection}{完全交}。则或者
\begin{enumerate}
\item[(a)] $H^{2n}(V,\mathbb Q_\ell(n))$ 中的每个\Term{algebraic-cohomology-class}{代数上同调类}都是 $\eta^n$ 的有理数倍；
\item[(b)] $H^{2n}(V,\mathbb Q_\ell(n))$ 由\Term{algebraic-cohomology-class}{代数上同调类}生成。
\end{enumerate}
若（b）成立，则对\Term{characteristic}{特征}为 $p$ 的\Term{algebraically-closed-field}{代数闭域}上每个\Term{dimension}{维数}为 $2n$、\Term{multidegree}{多重次数}为 $\underline a$ 的\Term{complete-intersection}{完全交}，$H^{2n}(V,\mathbb Q_\ell(n))$ 都由\Term{algebraic-cohomology-class}{代数类}生成。
\end{SGAresult}

取 $X$ 为\Term{dimension}{维数} $2n+1$、\Term{multidegree}{多重次数} $(a_1\ldots a_{d-1})$ 的\Term{generic}{一般}\Term{complete-intersection}{完全交}，取 $\ell$ 为 $X$ 与\Term{degree}{次数}为 $a_d$ 的\Term{hypersurface}{超曲面}所成\Term{section}{截面}的\Term{generic}{一般}\Term{pencil}{铅笔}，应用1.4（参见第 XVII 讲2.4）。事实上，也可以取 $X$ 的\Term{multidegree}{多重次数}为 $(a_1\ldots a_d)$，而 $\ell$ 为\Term{hyperplane-section}{超平面截面}的\Term{generic}{一般}\Term{pencil}{铅笔}。最后的断言来自这样一个事实：\Term{algebraic-cohomology-class}{代数上同调类}经\Term{specialization}{专化}后仍是\Term{algebraic-cohomology-class}{代数的}。

\SGAsection{2}{对第 XVIII 讲的补充}
\begin{SGAresult}{命题}{2.1}
设 $X\subset\mathbb P$ 是\Term{algebraically-closed-field}{代数闭域} $k$ 上纯偶\Term{dimension}{维数} $2n$ 的\Term{smooth}{光滑}\Term{projective-variety}{射影簇}，$\ell\subset\check{\mathbb P}$ 是 $X$ 的\Term{hyperplane-section}{超平面截面}的\Term{generic}{一般}\Term{lefschetz-pencil}{Lefschetz 铅笔}。设 $s$ 是 $\ell$ 的一个\Term{geometric-point}{几何点}，$H_{s}$ 是相应的\Term{hyperplane-section}{超平面截面}，$S\subset\ell$ 是\Term{pencil}{铅笔}的\Term{exceptional-value}{例外值}\Term{set}{集合}，$\pi_1(\ell-S,s)$ 是\Term{monodromy-group}{单值群}。

假设 $X$ 满足（LV）。则 $\pi_1(\ell-S,s)$ 在\Term{vanishing-cohomology}{消失上同调}
\[
\operatorname{Ev}(H_{s},\mathbb Q_\ell)\subset H^{2n}(H_{s},\mathbb Q_\ell)
\]
上的\Term{representation}{表示}是\Term{absolutely-irreducible}{绝对不可约的}。
\end{SGAresult}
\begin{proof}
回顾 $\operatorname{Ev}(H_{s},\mathbb Q_\ell)$ 是 $H^{2n}(X,\mathbb Q_\ell)$ 在 $H^{2n}(H_{s},\mathbb Q_\ell)$ 中的像的\Term{orthogonal-complement}{正交补}；这个像等于\Term{monodromy-group}{单值群}\Term{action}{作用}下的\Term{invariant-subgroup}{不变子群}（第 XVIII 讲5.6.10）。

\SourcePage{342}
\Term{projective-line}{射影直线} $\mathbb P^1$ 是\Term{simply-connected}{单连通的}。因此 $\pi_1(\ell-S,s)$ 由 $S$ 各点的\Term{inertia-group}{惯性子群}及其\Term{conjugacy}{共轭}（在\Term{topology}{拓扑}意义下）生成。故其在 $\operatorname{GL}(H^n(H_{s},\mathbb Q_\ell))$ 中的像由变换
\[
x\longmapsto x-(-1)^n(x,\delta)\delta
\]
（在\Term{topology}{拓扑}意义下）生成，其中 $\delta$ 与某个\Term{vanishing-cycles}{消失圈}\Term{conjugacy}{共轭}。因此，$H^{2n}(X,\mathbb Q_\ell)$ 在 $H^{2n}(H_{s},\mathbb Q_\ell)$ 中的像，就是\Term{vanishing-cycles}{消失圈}及其\Term{conjugacy}{共轭}所生成\Term{vector-subspace}{子空间}的\Term{orthogonal-complement}{正交补}：$\operatorname{Ev}(H_{s},\mathbb Q_\ell)$ 由\Term{vanishing-cycles}{消失圈}及其\Term{conjugacy}{共轭}生成。

对于\Term{generic}{一般}\Term{lefschetz-pencil}{Lefschetz 铅笔}，第 XVIII 讲第6节的论证仍说明，对应于 $S$ 中不同点的\Term{vanishing-cycles}{消失圈} $\delta$ 彼此\Term{conjugacy}{共轭}。

设 $V\subset\operatorname{Ev}(H_{s},\mathbb Q_\ell)$ 是非零\Term{invariant-subspace}{不变子空间}。存在一个\Term{vanishing-cycles}{消失圈}的\Term{conjugacy}{共轭} $\delta$ 和 $v\in V$，使 $(v,\delta)\ne0$（这里用到了（LV））。由于 $v-(-1)^m(v,\delta)\delta\in V$，故 $\delta\in V$。\Term{vector-subspace}{子空间} $V$ 包含一个\Term{vanishing-cycles}{消失圈}（因而也包含其全部\Term{conjugacy}{共轭}），所以等于 $\operatorname{Ev}(H_{s},\mathbb Q_\ell)$。

\Term{scalar-coefficient-field}{系数域} $\mathbb Q_\ell$ \Term{group-extension}{扩张}后，这个论证仍然适用，从而得到\Term{absolute-irreducibility}{绝对不可约性}。
\end{proof}

\textbf{注记 2.2。}若去掉假设（LV），上述论证仍证明 $H^{2n}(H_{s},\mathbb Q_\ell)$ 的每个\Term{invariant-subspace}{不变子空间}或者包含 $\operatorname{Ev}(H_{s},\mathbb Q_\ell)$，或者包含在 $H^{2n}(X,\mathbb Q_\ell)$ 的像中。当 $X$ 为偶\Term{dimension}{维数}时，对任意\Term{lefschetz-pencil}{Lefschetz 铅笔}也有同样的结论。

\SGAsection{3}{1.3的证明}
\textbf{3.1。}所需证明的是，在\Term{theorem}{定理}的假设下1.5（b）不成立。$h^{0,2n}\ne0$ 的情形将由 Katz 在第 XXI 讲4.1中处理（条件 $h^{0,2n}\ne0$ \Term{equivalence}{等价}于 $\sum d_i-(2n+r+1)\geq0$，这里的数正是满足 $\Omega_V^{2n}\simeq\mathcal O(m)$ 的整数 $m$）。本讲的结果在研讨班后面不再使用，因而不存在循环论证。

\SourcePage{343}
\textbf{3.2。}当 $k=\mathbb C$ 时，只须注意：若1.5（b）成立，则 $V$ 的\Term{cohomology}{上同调}纯为 $(n,n)$ 型，因为\Term{algebraic-cohomology}{代数上同调}为 $(n,n)$ 型。这只在1.3中排除的例外情形才会发生（第 XI 讲2.9）。由\Term{lefschetz-principle}{Lefschetz 原理}，1.3在\Term{characteristic}{特征} $0$ 时成立。

\begin{SGAresult}{命题}{3.3}
在1.4的假设下，假定（b）成立。设 $\operatorname{Ev}(V,\mathbb Q)$ 是 $\operatorname{Ev}(V,\mathbb Q_\ell(n))$ 中由\Term{algebraic-cohomology-class}{代数上同调类}生成的 $\mathbb Q$-\Term{vector-subspace}{向量子空间}。
\begin{enumerate}
\item[(a)] 有
\[
\operatorname{Ev}(V,\mathbb Q)\otimes\mathbb Q_\ell
\xrightarrow{\sim}\operatorname{Ev}(V,\mathbb Q_\ell(n)).
\]
在 $\operatorname{Ev}(V,\mathbb Q)$ 上，\Term{intersection-form}{交形式}是有理的\Term{definite-form}{定号形式}，其\Term{signature}{符号}为 $(-1)^n$。
\item[(b)] \Term{vanishing-cycles}{消失圈}及其\Term{conjugacy}{共轭}属于 $\operatorname{Ev}(V,\mathbb Q)$，并在其中构成一个 $A$、$D$ 或 $E$ 型的\Term{irreducible-root-system}{不可约根系}。
\item[(c)] \Term{monodromy-group}{单值群}（即 $\pi_1(\ell-S,\bar s)$ 在 $\operatorname{GL}(\operatorname{Ev}(V,\mathbb Q))$ 中的像）是\Term{finite-group}{有限群}；它就是（b）中\Term{root-system}{根系}的\Term{weyl-group}{Weyl 群}。
\end{enumerate}
\end{SGAresult}
\begin{proof}
除\Term{intersection-form}{交形式}的\Term{definite-form}{定号性}外，（a）来自\Term{intersection-form}{交形式}在 $\operatorname{Ev}(V,\mathbb Q_\ell(n))$ 上\Term{nondegenerate-form}{非退化}（由（LV）），并且在 $\operatorname{Ev}(V,\mathbb Q)$ 上取有理值。

条件1.4（b）\Term{equivalence}{等价}于说：$V$ 上存在模\Term{numerical-equivalence}{数值等价}后不在\Term{monodromy}{单值性}\Term{action}{作用}下\Term{invariant}{不变}的\Term{algebraic-cycle}{代数圈}。因此这个条件与 $\ell$ 无关。取这样的一个\Term{algebraic-cycle}{代数圈} $x$；设 $\delta$ 是对应于 $s_0\in S$ 以及从定义 $V$ 的 $\bar s$ 到 $s_0$ 的路径 $\mathrm{ch}$ 的\Term{vanishing-cycles}{消失圈}。适当选取 $\mathrm{ch}$，可使 $(x,\delta)\ne0$。$x$ 的 Galois 变换
\[
\sigma(x)=x-(-1)^n(x,\delta)\delta
\]
仍是\Term{algebraic-cycle}{代数的}；故 $(x,\delta)\delta$ 是\Term{algebraic-cycle}{代数的}，是一个与 $\ell$ 无关的圈 $y$ 的像。因此有理数
\[
\frac{(-1)^n(y,y)}{2}=(x,\delta)^2
\]
与 $\ell$ 无关；它在所有 $\mathbb Q_\ell$ 中都是平方，因而是一个平方，并且 $\pm(x,\delta)$ 属于 $\mathbb Q$，且与 $\ell$ 无关。这不仅说明 $\delta\in\operatorname{Ev}(V,\mathbb Q)$，还说明模\Term{numerical-equivalence}{数值等价}的\Term{algebraic-cycle}{代数圈}
\[
\pm\delta=\frac{\sigma(x)-x}{(x,\delta)}
\]
与 $\ell$ 无关。

\SourcePage{344}
\Term{variety}{簇} $V$ 上的\Term{algebraic-cycle}{代数圈}总是在 $V$ 的任一\Term{field-of-definition}{定义域}的某个\Term{finite-extension}{有限扩张}上有定义。因此 $\pi_1(\ell-S,s)$ 经由一个\Term{finite-group}{有限群}\Term{action}{作用}在 $\operatorname{Ev}(V,\mathbb Q)$ 上。这个群由与\Term{vanishing-cycles}{消失圈} $\delta$ 及其\Term{conjugacy}{共轭}相联系的\Term{reflection}{反射}生成（第 XVIII 讲6.6.1及本讲第2节）。所以这些 $\delta$ 及其\Term{conjugacy}{共轭}构成一个\Term{root-system}{根系}，其\Term{weyl-group}{Weyl 群}就是\Term{monodromy-group}{单值群}。此外，这个\Term{root-system}{根系}的所有\Term{root}{根}在\Term{weyl-group}{Weyl 群}下彼此\Term{conjugacy}{共轭}，故它是 $A$、$D$ 或 $E$ 型的。

在相差一个因子的意义下，\Term{intersection-form}{交形式}是唯一在\Term{weyl-group}{Weyl 群}下\Term{invariant}{不变}的\Term{form}{形式}。因此它是\Term{definite-form}{定号的}。由
\[
(\delta,\delta)=(-1)^n\cdot2
\]
确定其\Term{signature}{符号}。
\end{proof}

\begin{SGAresult}{命题}{3.4}
若 $k=\mathbb C$，则在1.4的假设下，下列条件\Term{equivalence}{等价}：
\begin{enumerate}
\item[(i)] \Term{monodromy-group}{单值群}是\Term{finite-group}{有限群}；
\item[(ii)] \Term{vanishing-cohomology}{消失上同调} $\operatorname{Ev}(V)$ 纯为 $(n,n)$ 型。
\end{enumerate}
\end{SGAresult}
\begin{proof}
设 $\widetilde X\to\ell$ 是\Term{fibre}{纤维}为\Term{pencil}{铅笔} $\ell$ 中\Term{hyperplane-section}{超平面截面}的\Term{fibre-bundle}{纤维丛}，$V_s$ 是 $s\in\ell(\mathbb C)$ 上的\Term{fibre}{纤维}。各 $\operatorname{Ev}(V_s)$ 在 $\ell$ 上形成\Term{hodge-structure}{Hodge 结构}的\Term{algebraic-variation-of-hodge-structures}{代数变分}。将\Term{fixed-part-theorem}{固定部分定理} [2] 4.1.2应用于 $\ell$ 的一个\Term{finite-covering}{有限覆盖}（参见 [2] 4.1.3.3），可知（i）蕴含 $\operatorname{Ev}(V_s)$ 的\Term{hodge-structure}{Hodge 结构}是\Term{locally-constant}{局部常值的}。因而\Term{monodromy-group}{单值群}通过\Term{hodge-structure}{Hodge 结构}的\Term{automorphism}{自同构}\Term{action}{作用}。对一个\Term{picard-lefschetz-transformation}{Picard--Lefschetz 变换} $T$，$\operatorname{Im}(T-1)$ 是\Term{rank}{秩}为一的\Term{sub-hodge-structure}{子 Hodge 结构}，故为 $(n,n)$ 型，并由一个\Term{vanishing-cycles}{消失圈} $\delta$ 生成。既然\Term{vanishing-cycles}{消失圈}均为 $(n,n)$ 型，整个 $\operatorname{Ev}(V_s)$ 也是如此。

反过来，如果 $\operatorname{Ev}(V)$ 纯为 $(n,n)$ 型，则 $\operatorname{Ev}(V)$ 上的\Term{intersection-form}{交形式}是\Term{definite-form}{定号的}，而\Term{monodromy-group}{单值群}是\Term{finite-group}{有限群}（参见 [2] 4.1.3.3）。
\end{proof}

\SourcePage{345}
\textbf{注记 3.5。}若3.3中的\Term{equivalence}{等价}条件成立，则如3.3中一样可知，\Term{vanishing-cycles}{消失圈}（及其\Term{conjugacy}{共轭}）构成一个 $A$、$D$ 或 $E$ 型的\Term{root-system}{根系}，其\Term{monodromy-group}{单值群}就是\Term{weyl-group}{Weyl 群}。

\textbf{3.6。}下面对\Term{finite}{有限}\Term{characteristic}{特征} $\ne2$ 的 $k$ 证明1.3。设 $S=\operatorname{Spec}(\mathcal V)$ 是一个\Term{residue-field}{剩余域}为 $k$ 的\Term{mixed-characteristic}{混合特征}\Term{discrete-valuation-ring}{离散赋值环}的\Term{spectrum}{谱}，$X\subset\mathbb P_S^{2n+d}$ 是 $S$ 上\Term{relative-dimension}{相对维数}为 $2n+1$、\Term{multidegree}{多重次数}为 $(a_1\ldots a_{d-1})$ 的\Term{complete-intersection}{完全交}。设 $\ell$ 是 $X$ 的\Term{degree}{次数}为 $a_d$ 的\Term{hypersurface-section}{超曲面截面}的一个\Term{lefschetz-pencil}{Lefschetz 铅笔}（在\Term{special-fibre}{特殊纤维}上，$\ell_s$ 是 $X_s$ 的一个\Term{lefschetz-pencil}{Lefschetz 铅笔}）。设
\[
f:\widetilde X\longrightarrow\ell\simeq\mathbb P_S^1
\]
是\Term{fibre}{纤维}为\Term{pencil}{铅笔}中各\Term{section}{截面}的\Term{fibre-bundle}{纤维丛}。\Term{exceptional-locus}{例外轨迹} $T\subset\ell$ 在 $S$ 上是\Term{etale}{平展的}（第 XVII 讲第6节；这里用到 $p\ne2$）。\Term{sheaf}{层} $R^{2n}f_*\mathbb Q_\ell$ 在 $\ell-T$ 上\Term{locally-constant}{局部常值}；因此它是\Term{tamely-ramified}{温和分歧的}，其\Term{monodromy}{单值性}在\Term{characteristic}{特征} $p$ 与\Term{characteristic}{特征} $0$ 下相同。如果1.5（b）成立，则由3.3（c）和3.4可知，\Term{dimension}{维数}为 $2n$、\Term{multidegree}{多重次数}为 $\underline a$ 的复\Term{complete-intersection}{完全交}的\Term{cohomology-group}{上同调群} $H^{2n}(V)$ 纯为 $(n,n)$ 型，而这只在1.3中排除的那些情形才会发生（第 XI 讲2.9）。

\SGAsection{4}{对第 XVIII 讲的第二个补充}
本节在 $\mathbb C$ 上讨论，使用整系数（\Term{transcendental-cohomology}{超越}）\Term{cohomology}{上同调}。我们用更接近 Lefschetz 的方法重新推导第 XVIII 讲的某些结果。在第 XVIII 讲中，这些结果只在与 $\mathbb Q$（更确切地说，与 $\mathbb Q_\ell$，两者在这里是一样的）作\Term{tensor-product}{张量积}后得到了证明。

\textbf{4.1。}设 $X\subset\mathbb P$ 是纯\Term{dimension}{维数} $n+1$ 的\Term{smooth}{光滑}复\Term{projective-variety}{射影簇}，$\ell$ 是 $X$ 的\Term{hyperplane-section}{超平面截面}的一个\Term{lefschetz-pencil}{Lefschetz 铅笔}，其\Term{pencil-axis}{轴}为 $A$。也用 $\ell$ 表示此\Term{pencil}{铅笔}的\Term{parameter-space}{参数空间}，它与 $\mathbb P^1$ \Term{isomorphism}{同构}。设 $S\subset\ell$ 是\Term{exceptional-value}{例外值}\Term{set}{集合}。对 $s\in\ell$，以 $H_s$ 表示 $X$ 的相应\Term{hyperplane-section}{超平面截面}（其\Term{dimension}{维数}为 $n$）。

\SourcePage{346}
设 $T(A)$ 是 $A\cap X$ 在 $X$ 中的一个闭\Term{tubular-neighbourhood}{管状邻域}，其\Term{boundary}{边界}与所有 $H_s$ \Term{transverse}{横截}。

设 $0,\infty$ 是 $\ell-S$ 中的两个点，并设 $\alpha$ 是从 $\ell-S$ 到 $\mathbb P^1(\mathbb C)$ 的一个\Term{orientation-preserving-diffeomorphism}{保定向微分同胚}，将 $0$ 映到 $0$，$\infty$ 映到 $\infty$，并将 $S$ 映到以 $0$ 为\Term{central}{中心的}\Term{unit}{单位}圆周上。$\alpha$ 与到 $\ell$ 的\Term{projection}{投影}的复合 $\alpha f$，使 $X-T(A)^\circ$ 成为 $\mathbb P^1(\mathbb C)$ 上的一个\Term{manifold-with-boundary}{带边流形}的\Term{singular-fibration}{奇异纤维丛}（每个 $\alpha(S)$ 中的点上方都有一个\Term{singularity}{奇点}）。以 $c_s$ 表示连接 $0$ 与 $\alpha(s)$ 的线段。

省略 $\alpha$ 不写，图示如下：
\begin{center}
\begin{tikzpicture}[scale=1.15]
\coordinate (O) at (0,0);
\draw (O)--(1.65,1.4) node[right] {$s_1$};
\draw (O)--(0,2) node[above] {$s_2$};
\draw (O)--(-1.25,1.4);
\node[below right] at (O) {$0$};
\node at (1.15,.62) {$c_1$};
\node[right] at (0,1.05) {$c_2$};
\node at (-.36,.69) {$\cdots$};
\end{tikzpicture}
\end{center}

在 $s\in S$ 及 $H_s\subset X$ 的\Term{singularity}{奇点} $x_s$ 附近，适当选取\Term{local-coordinates}{局部坐标}后，就处在第 XIV 讲3.2所研究的情形中。若选取的 $\alpha$ 在 $s$ 附近\Term{holomorphic}{全纯}，甚至可以在 $X$ 与 $\ell$ 上找到\Term{local-coordinates}{局部坐标}，使 $f$ 写为 $\sum z_i^2$，而 $\alpha^{-1}(c_s)$ 由“$z$ 为正实数”给出。在 $s$ 的某个\Term{neighborhood}{邻域}与 $\alpha^{-1}(c_s)-s$ 的交集上方，可以考虑由“$z_i$ 为实数”定义的实\Term{sphere-bundle}{球面丛} $\Delta'_s\subset X$。当 $t\to s$ 时，$\Delta'_s$ 在 $t$ 处的\Term{fibre}{纤维} $\Delta'_{st}$ 趋于 $x_s$。以 $\Delta_s$ 表示 $\{x_s\}$ 与 $\alpha^{-1}(c_s)-\{s\}$ 上一个\Term{sphere-bundle}{球面丛}的并，后者包含在 $X-T(A)$ 中并延拓 $\Delta'_s$。$\Delta_s$ 是\Term{embedding}{嵌入} $X$ 中的一个\Term{disk}{圆盘}，其\Term{boundary}{边界} $H_0\cap\Delta_s$ 是一个\Term{sphere}{球面}，代表 $s$ 处的\Term{vanishing-cycles}{消失圈} $\delta_s$ 沿路径 $\alpha^{-1}(c_s)$ \Term{parallel-transport}{平移}到 $H_0$ 中所得的圈。

对函数 $|\alpha f|^2$ 应用\Term{morse-theory}{Morse 理论}，以研究当 $r$ 增大时
\[
\{x\in X\mid x\in T(A)\ \text{或}\ |\alpha f(x)|^2\leq r\}
\]
的\Term{homotopy}{同伦}如何变化。

\SourcePage{347}
令
\[
X^- =\{x\in X\mid x\in T(A)\ \text{或}\ \alpha f(x)\ne\infty\}.
\]
可得\Term{inclusion}{包含映射}
\[
H_0\cup\bigcup_{s\in S}\Delta_s
\lhook\joinrel\longrightarrow
T(A)\cup H_0\cup\bigcup_{s\in S}\Delta_s
\lhook\joinrel\longrightarrow X^-
\]
是\Term{homotopy-equivalence}{同伦等价}。又由\Term{morse-theory}{Morse 理论}可知，$X-X^-$ 可以\Term{deformation-retraction}{形变收缩}到一个实\Term{dimension}{维数} $\leq n$ 的\Term{subcomplex}{子复形}上。综上，得到下述结果。

\begin{SGAresult}{\Term{theorem}{定理}}{4.2}
在 $H_\infty^{\mathrm{aff}}=H_\infty-A$ 中存在一个\Term{compact}{紧}的、（实）\Term{dimension}{维数} $\leq n$ 的\Term{subcomplex}{子复形} $\Sigma$，使\Term{inclusion}{包含映射}
\[
H_0\cup\bigcup_{s\in S}\Delta_s\longrightarrow X-\Sigma
\]
为\Term{homotopy-equivalence}{同伦等价}。
\end{SGAresult}

\textbf{4.3。}去掉 $\Sigma$ 不影响 $X$ 的指标 $i\leq n$ 的\Term{cohomology-group}{上同调群}、\Term{homology-group}{同调群}或\Term{homotopy-group}{同伦群}。由于 $H_0\cup\bigcup_{s\in S}\Delta_s$ 是由 $H_0$ 沿代表\Term{vanishing-cycles}{消失圈}的\Term{sphere}{球面}粘接若干 $(n+1)$ 维\Term{disk}{圆盘}得到的，在\Term{homology}{同调}中得到
\[
\mathbb Z^S\longrightarrow H_n(V_0)\longrightarrow H_n(X)\longrightarrow0.
\]
\Term{cohomology}{上同调}中相应的\Term{exact-sequence}{正合列}，以及由它经\Term{poincare-duality}{Poincaré 对偶}得到的\Term{exact-sequence}{正合列}，组成如下\Term{commutative-diagram}{交换图}：
\begin{equation}
\begin{tikzcd}[column sep=large,row sep=large]
 & \mathbb Z^S & & \\
\mathbb Z^S\arrow[r,"{\sum n_s\delta_s}"] & H^n(V_0)\arrow[u,"{(\,\cdot\,,\delta_s)}"']\arrow[r] & H^{n+2}(X)\arrow[r] & 0\\
 & H^n(X)\arrow[u]\arrow[ur,"{\eta\wedge}"'] & & \\
 & 0\arrow[u] & &
\end{tikzcd}
\tag{4.3.1}
\end{equation}

\SourcePage{348}
\SGAsection{5}{Noether 定理的例外情形}
本节继续在 $\mathbb C$ 上讨论，并使用整系数\Term{cohomology}{上同调}。

\textbf{5.1。}设 $2n,d$ 和 $\underline a=(a_1,\ldots,a_d)$（$d\geq1$，$a_i\geq2$）使得\Term{dimension}{维数}为 $2n$、\Term{multidegree}{多重次数}为 $\underline a$ 的\Term{smooth}{光滑}\Term{complete-intersection}{完全交} $V$ 的 $2n$ 维\Term{cohomology}{上同调}纯为 $(n,n)$ 型。换言之，$n=0$，或者 $(2n;\underline a)$ 等于 $(2n;2)$、$(2n;2,2)$ 或 $(2;3)$（第 XI 讲2.9）。

假设 $V$ 由一个\Term{dimension}{维数}为 $2n+1$、\Term{multidegree}{多重次数}为 $(a_1,\ldots,a_{d-1})$ 的\Term{complete-intersection}{完全交} $X\subset\mathbb P^{2n+d}$（相应地，由\Term{multidegree}{多重次数}为 $(a_1,\ldots,a_d)$ 的\Term{complete-intersection}{完全交} $X\subset\mathbb P^{2n+d+1}$），取\Term{lefschetz-pencil}{Lefschetz 铅笔} $\ell$ 中一个\Term{degree}{次数}为 $a_d$ 的\Term{hypersurface-section}{超曲面截面}（相应地，取一个\Term{hyperplane-section}{超平面截面}）而得。

$H^{2n}(V)$ 的\Term{vanishing-cohomology}{消失部分} $\operatorname{Ev}(V)$ 等于此\Term{cohomology}{上同调}的\Term{primitive-part}{本原部分}；根据4.3.1，它在 $\mathbb Z$ 上由\Term{vanishing-cycles}{消失圈}生成。它是\Term{hyperplane-section}{超平面截面}的\Term{cohomology-class}{上同调类} $\eta$ 的 $n$ 次幂 $\eta^n$ 的\Term{orthogonal-complement}{正交补}。

由于 $H^{2n}(V)$ 上的\Term{intersection-form}{交形式}是\Term{unimodular}{幺模的}，且 $\eta^{2n}$ 不是另一个\Term{cohomology-class}{上同调类}的 $b$ 倍（$b>1$）（第 XI 讲1.6（iv）），故 $\operatorname{Ev}(V)$ 上\Term{intersection-form}{交形式}的\Term{discriminant}{判别式}等于
\[
(\eta^n,\eta^n)=\prod a_i.
\]

设 $R$ 是 $H^{2n}(V)$ 中由\Term{vanishing-cycles}{消失圈}及其\Term{conjugacy}{共轭}构成的\Term{root-system}{根系}。关于 $R$，我们知道：
\begin{enumerate}
\item[(a)] $R$ 为空，或为 $A$、$D$、$E$ 型；
\item[(b)] $R$ 的\Term{rank}{秩}已在第 XI 讲1.10中给出；
\item[(c)] \Term{discriminant}{判别式}为 $\prod a_i$。在表中，这个\Term{discriminant}{判别式}是相应的\Term{simply-connected}{单连通}复\Term{simple-group}{单群}的\Term{center-of-group}{中心}的阶。我们计算\Term{discriminant}{判别式}时的论证说明，这个\Term{center-of-group}{中心}是\Term{cyclic-group}{循环群}；它是\Term{weight-lattice}{权格}模\Term{root-lattice}{根格}之商的\Term{dual-group}{对偶群}。
\end{enumerate}

\SourcePage{349}
在所涉及的所有情形中，若 $n\ne0$，这些信息完全确定 $R$。可得：
\begin{SGAresult}{命题}{5.2}
\begin{enumerate}
\item[(i)] 若 $V$ 是偶\Term{dimension}{维数}的\Term{quadric}{二次超曲面}，则 $R$ 为 $A_1$ 型。
\item[(ii)] 若 $V$ 是\Term{cubic-surface}{三次曲面}，则 $R$ 为 $E_6$ 型。
\item[(iii)] 若 $V$ 的\Term{dimension}{维数}为 $2n$，并且是两个\Term{quadric}{二次超曲面}的交，则 $R$ 为 $D_{2n+3}$ 型。
\end{enumerate}
\end{SGAresult}

\textbf{5.3。}假设 $n\ne0$。$\operatorname{Ev}(V)$ 中\Term{length}{长度}为 $\pm2$ 的向量所成的\Term{set}{集合} $R^1$ 也是一个\Term{root-system}{根系}，它具有与 $R$ 相同的\Term{characteristic}{特征}（a）、（b）、（c）。因而 $R=R^1$。设 $\pi_1$ 是所讨论类型的所有\Term{complete-intersection}{完全交} $V\subset\mathbb P^{2n+d}$ 的\Term{parameter-space}{参数空间}的\Term{fundamental-group}{基本群}。$\pi_1$ 在 $\operatorname{GL}(\operatorname{Ev}(V))$ 中的像是 $R$ 的\Term{weyl-group}{Weyl 群}。在所涉及的所有情形中，它就是三元组
\[
\bigl(H^{2n}(V),\ \eta^n,\ \text{交形式}\bigr)
\]
的\Term{automorphism-group}{自同构群}。

\textbf{5.4。}将 $R$ 换成5.3中的 $R^1$ 后，4.2（ii）的结果是熟知的 [1]；而对较小的 $n$，4.2（iii）至少已为 Manin 所知（$n=0$ 的情形是显然的；两个\Term{quadric}{二次超曲面}相交所得的\Term{surface}{曲面}是一个\Term{del-pezzo-surface}{del Pezzo 曲面}）。M. Reid [4] 对两个\Term{quadric}{二次超曲面}的交作了详细研究。他通过研究两个\Term{quadric}{二次超曲面}的交 $V\subset\mathbb P^{2n+2}$ 上 $n$ 维\Term{linear-subspace}{线性子空间}的几何，推\Term{derive}{导出}4.2（iii）。

\textbf{参考文献}
\begin{enumerate}
\item[{[1]}] E. CARTAN，关于平面四次\Term{curve}{曲线}的28条双切线与\Term{cubic-surface}{三次曲面}的27条直线的几点注记（Quelques remarques sur les 28 bitangentes d'une quartique plane et les 27 droites d'une surface cubique），\textit{Bull. Sc. Math.} \textbf{70}（1946），第42--45页；\textit{Œuvres}，第一部分，第2卷，第1353页。
\item[{[2]}] P. DELIGNE，\Term{hodge-theory}{Hodge 理论} II（Théorie de Hodge II），\textit{Publ. Math. IHES} \textbf{40}（1971），第5--58页。
\end{enumerate}

\SourcePage{350}
\begin{enumerate}
\item[{[3]}] S. LEFSCHETZ，位置分析与\Term{algebraic-geometry}{代数几何}（L'analysis situs et la géométrie algébrique），Gauthier-Villars，1924（收入 \textit{Selected papers} 重印，Chelsea）。
\item[{[4]}] M. REID，学位论文，Cambridge，1972。
\item[{[5]}] J. P. SERRE，\Term{coherent-algebraic-sheaf}{凝聚代数层}（Faisceaux algébriques cohérents），\textit{Ann. of Math.} \textbf{LXI}（1955），第197--278页（文中简称 FAC）。
\end{enumerate}
